% ===========================================================================
%  Shared preamble for the Interproof demo.
%
%  It sits OUTSIDE either document's root, which is the point: a real project
%  keeps one notation file and both the paper and the notes read it, so the
%  build has to cope with a source file that is not under `[[document]].root`.
%  SyncTeX records it as an input, Interproof drops it (it holds no labels and
%  no page of its own), and everything keeps working.
%
%  The two documents reach it two different ways, on purpose:
%    tex/paper/main.tex  % ===========================================================================
%  Shared preamble for the Interproof demo.
%
%  It sits OUTSIDE either document's root, which is the point: a real project
%  keeps one notation file and both the paper and the notes read it, so the
%  build has to cope with a source file that is not under `[[document]].root`.
%  SyncTeX records it as an input, Interproof drops it (it holds no labels and
%  no page of its own), and everything keeps working.
%
%  The two documents reach it two different ways, on purpose:
%    tex/paper/main.tex  % ===========================================================================
%  Shared preamble for the Interproof demo.
%
%  It sits OUTSIDE either document's root, which is the point: a real project
%  keeps one notation file and both the paper and the notes read it, so the
%  build has to cope with a source file that is not under `[[document]].root`.
%  SyncTeX records it as an input, Interproof drops it (it holds no labels and
%  no page of its own), and everything keeps working.
%
%  The two documents reach it two different ways, on purpose:
%    tex/paper/main.tex  % ===========================================================================
%  Shared preamble for the Interproof demo.
%
%  It sits OUTSIDE either document's root, which is the point: a real project
%  keeps one notation file and both the paper and the notes read it, so the
%  build has to cope with a source file that is not under `[[document]].root`.
%  SyncTeX records it as an input, Interproof drops it (it holds no labels and
%  no page of its own), and everything keeps working.
%
%  The two documents reach it two different ways, on purpose:
%    tex/paper/main.tex  \input{preamble}              + TEXINPUTS=../common:
%    tex/note/main.tex   \input{../common/preamble}    plain relative path
%  See the comments in ../../interproof.toml.
% ===========================================================================

\usepackage[T1]{fontenc}
\usepackage[margin=1.1in]{geometry}
\usepackage{amsmath,amssymb,amsthm}
\usepackage{hyperref}
\usepackage{cleveref}          % MUST be loaded after hyperref

\hypersetup{colorlinks=true, linkcolor=[rgb]{0.10,0.25,0.55},
            citecolor=[rgb]{0.10,0.25,0.55}}

% ---- theorem environments -------------------------------------------------
% One shared counter, so `Definition 2.1` and `Lemma 2.2` interleave the way a
% reader expects.  The environment names here are exactly the ones listed in
% `[grammar].environments`; adding one means adding it in both places.
\theoremstyle{definition}
\newtheorem{definition}{Definition}[section]
\newtheorem{example}[definition]{Example}
\theoremstyle{plain}
\newtheorem{theorem}[definition]{Theorem}
\newtheorem{lemma}[definition]{Lemma}
\newtheorem{proposition}[definition]{Proposition}
\newtheorem{corollary}[definition]{Corollary}
\theoremstyle{remark}
\newtheorem{remark}[definition]{Remark}

\crefname{definition}{Definition}{Definitions}
\Crefname{definition}{Definition}{Definitions}
\crefname{lemma}{Lemma}{Lemmas}
\Crefname{lemma}{Lemma}{Lemmas}
\crefname{proposition}{Proposition}{Propositions}
\Crefname{proposition}{Proposition}{Propositions}
\crefname{corollary}{Corollary}{Corollaries}
\Crefname{corollary}{Corollary}{Corollaries}
\crefname{remark}{Remark}{Remarks}
\Crefname{remark}{Remark}{Remarks}

% ---- notation -------------------------------------------------------------
\newcommand{\Var}{\mathsf{Var}}
\newcommand{\Nat}{\mathbb{N}}
\newcommand{\Stores}{\Sigma}
\newcommand{\Assn}{\mathcal{A}}
\newcommand{\kw}[1]{\mathsf{#1}}
\newcommand{\skipk}{\kw{skip}}
\newcommand{\ifk}{\kw{if}}
\newcommand{\thenk}{\kw{then}}
\newcommand{\elsek}{\kw{else}}
\newcommand{\whilek}{\kw{while}}
\newcommand{\dok}{\kw{do}}
\newcommand{\prock}{\kw{proc}}
\newcommand{\sem}[1]{[\![#1]\!]}
\newcommand{\upd}[3]{#1[#2 \mapsto #3]}          % sigma[x -> n]
\newcommand{\bigstep}[3]{\langle #1, #2\rangle \Downarrow #3}
\newcommand{\triple}[3]{\{#1\}\, #2 \,\{#3\}}
\newcommand{\vdashH}{\vdash_{\mathrm{H}}}        % derivable in the proof system
\newcommand{\agree}[3]{#2 \sim_{#1} #3}          % agreement on a frame
\newcommand{\fv}{\operatorname{fv}}
              + TEXINPUTS=../common:
%    tex/note/main.tex   % ===========================================================================
%  Shared preamble for the Interproof demo.
%
%  It sits OUTSIDE either document's root, which is the point: a real project
%  keeps one notation file and both the paper and the notes read it, so the
%  build has to cope with a source file that is not under `[[document]].root`.
%  SyncTeX records it as an input, Interproof drops it (it holds no labels and
%  no page of its own), and everything keeps working.
%
%  The two documents reach it two different ways, on purpose:
%    tex/paper/main.tex  \input{preamble}              + TEXINPUTS=../common:
%    tex/note/main.tex   \input{../common/preamble}    plain relative path
%  See the comments in ../../interproof.toml.
% ===========================================================================

\usepackage[T1]{fontenc}
\usepackage[margin=1.1in]{geometry}
\usepackage{amsmath,amssymb,amsthm}
\usepackage{hyperref}
\usepackage{cleveref}          % MUST be loaded after hyperref

\hypersetup{colorlinks=true, linkcolor=[rgb]{0.10,0.25,0.55},
            citecolor=[rgb]{0.10,0.25,0.55}}

% ---- theorem environments -------------------------------------------------
% One shared counter, so `Definition 2.1` and `Lemma 2.2` interleave the way a
% reader expects.  The environment names here are exactly the ones listed in
% `[grammar].environments`; adding one means adding it in both places.
\theoremstyle{definition}
\newtheorem{definition}{Definition}[section]
\newtheorem{example}[definition]{Example}
\theoremstyle{plain}
\newtheorem{theorem}[definition]{Theorem}
\newtheorem{lemma}[definition]{Lemma}
\newtheorem{proposition}[definition]{Proposition}
\newtheorem{corollary}[definition]{Corollary}
\theoremstyle{remark}
\newtheorem{remark}[definition]{Remark}

\crefname{definition}{Definition}{Definitions}
\Crefname{definition}{Definition}{Definitions}
\crefname{lemma}{Lemma}{Lemmas}
\Crefname{lemma}{Lemma}{Lemmas}
\crefname{proposition}{Proposition}{Propositions}
\Crefname{proposition}{Proposition}{Propositions}
\crefname{corollary}{Corollary}{Corollaries}
\Crefname{corollary}{Corollary}{Corollaries}
\crefname{remark}{Remark}{Remarks}
\Crefname{remark}{Remark}{Remarks}

% ---- notation -------------------------------------------------------------
\newcommand{\Var}{\mathsf{Var}}
\newcommand{\Nat}{\mathbb{N}}
\newcommand{\Stores}{\Sigma}
\newcommand{\Assn}{\mathcal{A}}
\newcommand{\kw}[1]{\mathsf{#1}}
\newcommand{\skipk}{\kw{skip}}
\newcommand{\ifk}{\kw{if}}
\newcommand{\thenk}{\kw{then}}
\newcommand{\elsek}{\kw{else}}
\newcommand{\whilek}{\kw{while}}
\newcommand{\dok}{\kw{do}}
\newcommand{\prock}{\kw{proc}}
\newcommand{\sem}[1]{[\![#1]\!]}
\newcommand{\upd}[3]{#1[#2 \mapsto #3]}          % sigma[x -> n]
\newcommand{\bigstep}[3]{\langle #1, #2\rangle \Downarrow #3}
\newcommand{\triple}[3]{\{#1\}\, #2 \,\{#3\}}
\newcommand{\vdashH}{\vdash_{\mathrm{H}}}        % derivable in the proof system
\newcommand{\agree}[3]{#2 \sim_{#1} #3}          % agreement on a frame
\newcommand{\fv}{\operatorname{fv}}
    plain relative path
%  See the comments in ../../interproof.toml.
% ===========================================================================

\usepackage[T1]{fontenc}
\usepackage[margin=1.1in]{geometry}
\usepackage{amsmath,amssymb,amsthm}
\usepackage{hyperref}
\usepackage{cleveref}          % MUST be loaded after hyperref

\hypersetup{colorlinks=true, linkcolor=[rgb]{0.10,0.25,0.55},
            citecolor=[rgb]{0.10,0.25,0.55}}

% ---- theorem environments -------------------------------------------------
% One shared counter, so `Definition 2.1` and `Lemma 2.2` interleave the way a
% reader expects.  The environment names here are exactly the ones listed in
% `[grammar].environments`; adding one means adding it in both places.
\theoremstyle{definition}
\newtheorem{definition}{Definition}[section]
\newtheorem{example}[definition]{Example}
\theoremstyle{plain}
\newtheorem{theorem}[definition]{Theorem}
\newtheorem{lemma}[definition]{Lemma}
\newtheorem{proposition}[definition]{Proposition}
\newtheorem{corollary}[definition]{Corollary}
\theoremstyle{remark}
\newtheorem{remark}[definition]{Remark}

\crefname{definition}{Definition}{Definitions}
\Crefname{definition}{Definition}{Definitions}
\crefname{lemma}{Lemma}{Lemmas}
\Crefname{lemma}{Lemma}{Lemmas}
\crefname{proposition}{Proposition}{Propositions}
\Crefname{proposition}{Proposition}{Propositions}
\crefname{corollary}{Corollary}{Corollaries}
\Crefname{corollary}{Corollary}{Corollaries}
\crefname{remark}{Remark}{Remarks}
\Crefname{remark}{Remark}{Remarks}

% ---- notation -------------------------------------------------------------
\newcommand{\Var}{\mathsf{Var}}
\newcommand{\Nat}{\mathbb{N}}
\newcommand{\Stores}{\Sigma}
\newcommand{\Assn}{\mathcal{A}}
\newcommand{\kw}[1]{\mathsf{#1}}
\newcommand{\skipk}{\kw{skip}}
\newcommand{\ifk}{\kw{if}}
\newcommand{\thenk}{\kw{then}}
\newcommand{\elsek}{\kw{else}}
\newcommand{\whilek}{\kw{while}}
\newcommand{\dok}{\kw{do}}
\newcommand{\prock}{\kw{proc}}
\newcommand{\sem}[1]{[\![#1]\!]}
\newcommand{\upd}[3]{#1[#2 \mapsto #3]}          % sigma[x -> n]
\newcommand{\bigstep}[3]{\langle #1, #2\rangle \Downarrow #3}
\newcommand{\triple}[3]{\{#1\}\, #2 \,\{#3\}}
\newcommand{\vdashH}{\vdash_{\mathrm{H}}}        % derivable in the proof system
\newcommand{\agree}[3]{#2 \sim_{#1} #3}          % agreement on a frame
\newcommand{\fv}{\operatorname{fv}}
              + TEXINPUTS=../common:
%    tex/note/main.tex   % ===========================================================================
%  Shared preamble for the Interproof demo.
%
%  It sits OUTSIDE either document's root, which is the point: a real project
%  keeps one notation file and both the paper and the notes read it, so the
%  build has to cope with a source file that is not under `[[document]].root`.
%  SyncTeX records it as an input, Interproof drops it (it holds no labels and
%  no page of its own), and everything keeps working.
%
%  The two documents reach it two different ways, on purpose:
%    tex/paper/main.tex  % ===========================================================================
%  Shared preamble for the Interproof demo.
%
%  It sits OUTSIDE either document's root, which is the point: a real project
%  keeps one notation file and both the paper and the notes read it, so the
%  build has to cope with a source file that is not under `[[document]].root`.
%  SyncTeX records it as an input, Interproof drops it (it holds no labels and
%  no page of its own), and everything keeps working.
%
%  The two documents reach it two different ways, on purpose:
%    tex/paper/main.tex  \input{preamble}              + TEXINPUTS=../common:
%    tex/note/main.tex   \input{../common/preamble}    plain relative path
%  See the comments in ../../interproof.toml.
% ===========================================================================

\usepackage[T1]{fontenc}
\usepackage[margin=1.1in]{geometry}
\usepackage{amsmath,amssymb,amsthm}
\usepackage{hyperref}
\usepackage{cleveref}          % MUST be loaded after hyperref

\hypersetup{colorlinks=true, linkcolor=[rgb]{0.10,0.25,0.55},
            citecolor=[rgb]{0.10,0.25,0.55}}

% ---- theorem environments -------------------------------------------------
% One shared counter, so `Definition 2.1` and `Lemma 2.2` interleave the way a
% reader expects.  The environment names here are exactly the ones listed in
% `[grammar].environments`; adding one means adding it in both places.
\theoremstyle{definition}
\newtheorem{definition}{Definition}[section]
\newtheorem{example}[definition]{Example}
\theoremstyle{plain}
\newtheorem{theorem}[definition]{Theorem}
\newtheorem{lemma}[definition]{Lemma}
\newtheorem{proposition}[definition]{Proposition}
\newtheorem{corollary}[definition]{Corollary}
\theoremstyle{remark}
\newtheorem{remark}[definition]{Remark}

\crefname{definition}{Definition}{Definitions}
\Crefname{definition}{Definition}{Definitions}
\crefname{lemma}{Lemma}{Lemmas}
\Crefname{lemma}{Lemma}{Lemmas}
\crefname{proposition}{Proposition}{Propositions}
\Crefname{proposition}{Proposition}{Propositions}
\crefname{corollary}{Corollary}{Corollaries}
\Crefname{corollary}{Corollary}{Corollaries}
\crefname{remark}{Remark}{Remarks}
\Crefname{remark}{Remark}{Remarks}

% ---- notation -------------------------------------------------------------
\newcommand{\Var}{\mathsf{Var}}
\newcommand{\Nat}{\mathbb{N}}
\newcommand{\Stores}{\Sigma}
\newcommand{\Assn}{\mathcal{A}}
\newcommand{\kw}[1]{\mathsf{#1}}
\newcommand{\skipk}{\kw{skip}}
\newcommand{\ifk}{\kw{if}}
\newcommand{\thenk}{\kw{then}}
\newcommand{\elsek}{\kw{else}}
\newcommand{\whilek}{\kw{while}}
\newcommand{\dok}{\kw{do}}
\newcommand{\prock}{\kw{proc}}
\newcommand{\sem}[1]{[\![#1]\!]}
\newcommand{\upd}[3]{#1[#2 \mapsto #3]}          % sigma[x -> n]
\newcommand{\bigstep}[3]{\langle #1, #2\rangle \Downarrow #3}
\newcommand{\triple}[3]{\{#1\}\, #2 \,\{#3\}}
\newcommand{\vdashH}{\vdash_{\mathrm{H}}}        % derivable in the proof system
\newcommand{\agree}[3]{#2 \sim_{#1} #3}          % agreement on a frame
\newcommand{\fv}{\operatorname{fv}}
              + TEXINPUTS=../common:
%    tex/note/main.tex   % ===========================================================================
%  Shared preamble for the Interproof demo.
%
%  It sits OUTSIDE either document's root, which is the point: a real project
%  keeps one notation file and both the paper and the notes read it, so the
%  build has to cope with a source file that is not under `[[document]].root`.
%  SyncTeX records it as an input, Interproof drops it (it holds no labels and
%  no page of its own), and everything keeps working.
%
%  The two documents reach it two different ways, on purpose:
%    tex/paper/main.tex  \input{preamble}              + TEXINPUTS=../common:
%    tex/note/main.tex   \input{../common/preamble}    plain relative path
%  See the comments in ../../interproof.toml.
% ===========================================================================

\usepackage[T1]{fontenc}
\usepackage[margin=1.1in]{geometry}
\usepackage{amsmath,amssymb,amsthm}
\usepackage{hyperref}
\usepackage{cleveref}          % MUST be loaded after hyperref

\hypersetup{colorlinks=true, linkcolor=[rgb]{0.10,0.25,0.55},
            citecolor=[rgb]{0.10,0.25,0.55}}

% ---- theorem environments -------------------------------------------------
% One shared counter, so `Definition 2.1` and `Lemma 2.2` interleave the way a
% reader expects.  The environment names here are exactly the ones listed in
% `[grammar].environments`; adding one means adding it in both places.
\theoremstyle{definition}
\newtheorem{definition}{Definition}[section]
\newtheorem{example}[definition]{Example}
\theoremstyle{plain}
\newtheorem{theorem}[definition]{Theorem}
\newtheorem{lemma}[definition]{Lemma}
\newtheorem{proposition}[definition]{Proposition}
\newtheorem{corollary}[definition]{Corollary}
\theoremstyle{remark}
\newtheorem{remark}[definition]{Remark}

\crefname{definition}{Definition}{Definitions}
\Crefname{definition}{Definition}{Definitions}
\crefname{lemma}{Lemma}{Lemmas}
\Crefname{lemma}{Lemma}{Lemmas}
\crefname{proposition}{Proposition}{Propositions}
\Crefname{proposition}{Proposition}{Propositions}
\crefname{corollary}{Corollary}{Corollaries}
\Crefname{corollary}{Corollary}{Corollaries}
\crefname{remark}{Remark}{Remarks}
\Crefname{remark}{Remark}{Remarks}

% ---- notation -------------------------------------------------------------
\newcommand{\Var}{\mathsf{Var}}
\newcommand{\Nat}{\mathbb{N}}
\newcommand{\Stores}{\Sigma}
\newcommand{\Assn}{\mathcal{A}}
\newcommand{\kw}[1]{\mathsf{#1}}
\newcommand{\skipk}{\kw{skip}}
\newcommand{\ifk}{\kw{if}}
\newcommand{\thenk}{\kw{then}}
\newcommand{\elsek}{\kw{else}}
\newcommand{\whilek}{\kw{while}}
\newcommand{\dok}{\kw{do}}
\newcommand{\prock}{\kw{proc}}
\newcommand{\sem}[1]{[\![#1]\!]}
\newcommand{\upd}[3]{#1[#2 \mapsto #3]}          % sigma[x -> n]
\newcommand{\bigstep}[3]{\langle #1, #2\rangle \Downarrow #3}
\newcommand{\triple}[3]{\{#1\}\, #2 \,\{#3\}}
\newcommand{\vdashH}{\vdash_{\mathrm{H}}}        % derivable in the proof system
\newcommand{\agree}[3]{#2 \sim_{#1} #3}          % agreement on a frame
\newcommand{\fv}{\operatorname{fv}}
    plain relative path
%  See the comments in ../../interproof.toml.
% ===========================================================================

\usepackage[T1]{fontenc}
\usepackage[margin=1.1in]{geometry}
\usepackage{amsmath,amssymb,amsthm}
\usepackage{hyperref}
\usepackage{cleveref}          % MUST be loaded after hyperref

\hypersetup{colorlinks=true, linkcolor=[rgb]{0.10,0.25,0.55},
            citecolor=[rgb]{0.10,0.25,0.55}}

% ---- theorem environments -------------------------------------------------
% One shared counter, so `Definition 2.1` and `Lemma 2.2` interleave the way a
% reader expects.  The environment names here are exactly the ones listed in
% `[grammar].environments`; adding one means adding it in both places.
\theoremstyle{definition}
\newtheorem{definition}{Definition}[section]
\newtheorem{example}[definition]{Example}
\theoremstyle{plain}
\newtheorem{theorem}[definition]{Theorem}
\newtheorem{lemma}[definition]{Lemma}
\newtheorem{proposition}[definition]{Proposition}
\newtheorem{corollary}[definition]{Corollary}
\theoremstyle{remark}
\newtheorem{remark}[definition]{Remark}

\crefname{definition}{Definition}{Definitions}
\Crefname{definition}{Definition}{Definitions}
\crefname{lemma}{Lemma}{Lemmas}
\Crefname{lemma}{Lemma}{Lemmas}
\crefname{proposition}{Proposition}{Propositions}
\Crefname{proposition}{Proposition}{Propositions}
\crefname{corollary}{Corollary}{Corollaries}
\Crefname{corollary}{Corollary}{Corollaries}
\crefname{remark}{Remark}{Remarks}
\Crefname{remark}{Remark}{Remarks}

% ---- notation -------------------------------------------------------------
\newcommand{\Var}{\mathsf{Var}}
\newcommand{\Nat}{\mathbb{N}}
\newcommand{\Stores}{\Sigma}
\newcommand{\Assn}{\mathcal{A}}
\newcommand{\kw}[1]{\mathsf{#1}}
\newcommand{\skipk}{\kw{skip}}
\newcommand{\ifk}{\kw{if}}
\newcommand{\thenk}{\kw{then}}
\newcommand{\elsek}{\kw{else}}
\newcommand{\whilek}{\kw{while}}
\newcommand{\dok}{\kw{do}}
\newcommand{\prock}{\kw{proc}}
\newcommand{\sem}[1]{[\![#1]\!]}
\newcommand{\upd}[3]{#1[#2 \mapsto #3]}          % sigma[x -> n]
\newcommand{\bigstep}[3]{\langle #1, #2\rangle \Downarrow #3}
\newcommand{\triple}[3]{\{#1\}\, #2 \,\{#3\}}
\newcommand{\vdashH}{\vdash_{\mathrm{H}}}        % derivable in the proof system
\newcommand{\agree}[3]{#2 \sim_{#1} #3}          % agreement on a frame
\newcommand{\fv}{\operatorname{fv}}
    plain relative path
%  See the comments in ../../interproof.toml.
% ===========================================================================

\usepackage[T1]{fontenc}
\usepackage[margin=1.1in]{geometry}
\usepackage{amsmath,amssymb,amsthm}
\usepackage{hyperref}
\usepackage{cleveref}          % MUST be loaded after hyperref

\hypersetup{colorlinks=true, linkcolor=[rgb]{0.10,0.25,0.55},
            citecolor=[rgb]{0.10,0.25,0.55}}

% ---- theorem environments -------------------------------------------------
% One shared counter, so `Definition 2.1` and `Lemma 2.2` interleave the way a
% reader expects.  The environment names here are exactly the ones listed in
% `[grammar].environments`; adding one means adding it in both places.
\theoremstyle{definition}
\newtheorem{definition}{Definition}[section]
\newtheorem{example}[definition]{Example}
\theoremstyle{plain}
\newtheorem{theorem}[definition]{Theorem}
\newtheorem{lemma}[definition]{Lemma}
\newtheorem{proposition}[definition]{Proposition}
\newtheorem{corollary}[definition]{Corollary}
\theoremstyle{remark}
\newtheorem{remark}[definition]{Remark}

\crefname{definition}{Definition}{Definitions}
\Crefname{definition}{Definition}{Definitions}
\crefname{lemma}{Lemma}{Lemmas}
\Crefname{lemma}{Lemma}{Lemmas}
\crefname{proposition}{Proposition}{Propositions}
\Crefname{proposition}{Proposition}{Propositions}
\crefname{corollary}{Corollary}{Corollaries}
\Crefname{corollary}{Corollary}{Corollaries}
\crefname{remark}{Remark}{Remarks}
\Crefname{remark}{Remark}{Remarks}

% ---- notation -------------------------------------------------------------
\newcommand{\Var}{\mathsf{Var}}
\newcommand{\Nat}{\mathbb{N}}
\newcommand{\Stores}{\Sigma}
\newcommand{\Assn}{\mathcal{A}}
\newcommand{\kw}[1]{\mathsf{#1}}
\newcommand{\skipk}{\kw{skip}}
\newcommand{\ifk}{\kw{if}}
\newcommand{\thenk}{\kw{then}}
\newcommand{\elsek}{\kw{else}}
\newcommand{\whilek}{\kw{while}}
\newcommand{\dok}{\kw{do}}
\newcommand{\prock}{\kw{proc}}
\newcommand{\sem}[1]{[\![#1]\!]}
\newcommand{\upd}[3]{#1[#2 \mapsto #3]}          % sigma[x -> n]
\newcommand{\bigstep}[3]{\langle #1, #2\rangle \Downarrow #3}
\newcommand{\triple}[3]{\{#1\}\, #2 \,\{#3\}}
\newcommand{\vdashH}{\vdash_{\mathrm{H}}}        % derivable in the proof system
\newcommand{\agree}[3]{#2 \sim_{#1} #3}          % agreement on a frame
\newcommand{\fv}{\operatorname{fv}}
              + TEXINPUTS=../common:
%    tex/note/main.tex   % ===========================================================================
%  Shared preamble for the Interproof demo.
%
%  It sits OUTSIDE either document's root, which is the point: a real project
%  keeps one notation file and both the paper and the notes read it, so the
%  build has to cope with a source file that is not under `[[document]].root`.
%  SyncTeX records it as an input, Interproof drops it (it holds no labels and
%  no page of its own), and everything keeps working.
%
%  The two documents reach it two different ways, on purpose:
%    tex/paper/main.tex  % ===========================================================================
%  Shared preamble for the Interproof demo.
%
%  It sits OUTSIDE either document's root, which is the point: a real project
%  keeps one notation file and both the paper and the notes read it, so the
%  build has to cope with a source file that is not under `[[document]].root`.
%  SyncTeX records it as an input, Interproof drops it (it holds no labels and
%  no page of its own), and everything keeps working.
%
%  The two documents reach it two different ways, on purpose:
%    tex/paper/main.tex  % ===========================================================================
%  Shared preamble for the Interproof demo.
%
%  It sits OUTSIDE either document's root, which is the point: a real project
%  keeps one notation file and both the paper and the notes read it, so the
%  build has to cope with a source file that is not under `[[document]].root`.
%  SyncTeX records it as an input, Interproof drops it (it holds no labels and
%  no page of its own), and everything keeps working.
%
%  The two documents reach it two different ways, on purpose:
%    tex/paper/main.tex  \input{preamble}              + TEXINPUTS=../common:
%    tex/note/main.tex   \input{../common/preamble}    plain relative path
%  See the comments in ../../interproof.toml.
% ===========================================================================

\usepackage[T1]{fontenc}
\usepackage[margin=1.1in]{geometry}
\usepackage{amsmath,amssymb,amsthm}
\usepackage{hyperref}
\usepackage{cleveref}          % MUST be loaded after hyperref

\hypersetup{colorlinks=true, linkcolor=[rgb]{0.10,0.25,0.55},
            citecolor=[rgb]{0.10,0.25,0.55}}

% ---- theorem environments -------------------------------------------------
% One shared counter, so `Definition 2.1` and `Lemma 2.2` interleave the way a
% reader expects.  The environment names here are exactly the ones listed in
% `[grammar].environments`; adding one means adding it in both places.
\theoremstyle{definition}
\newtheorem{definition}{Definition}[section]
\newtheorem{example}[definition]{Example}
\theoremstyle{plain}
\newtheorem{theorem}[definition]{Theorem}
\newtheorem{lemma}[definition]{Lemma}
\newtheorem{proposition}[definition]{Proposition}
\newtheorem{corollary}[definition]{Corollary}
\theoremstyle{remark}
\newtheorem{remark}[definition]{Remark}

\crefname{definition}{Definition}{Definitions}
\Crefname{definition}{Definition}{Definitions}
\crefname{lemma}{Lemma}{Lemmas}
\Crefname{lemma}{Lemma}{Lemmas}
\crefname{proposition}{Proposition}{Propositions}
\Crefname{proposition}{Proposition}{Propositions}
\crefname{corollary}{Corollary}{Corollaries}
\Crefname{corollary}{Corollary}{Corollaries}
\crefname{remark}{Remark}{Remarks}
\Crefname{remark}{Remark}{Remarks}

% ---- notation -------------------------------------------------------------
\newcommand{\Var}{\mathsf{Var}}
\newcommand{\Nat}{\mathbb{N}}
\newcommand{\Stores}{\Sigma}
\newcommand{\Assn}{\mathcal{A}}
\newcommand{\kw}[1]{\mathsf{#1}}
\newcommand{\skipk}{\kw{skip}}
\newcommand{\ifk}{\kw{if}}
\newcommand{\thenk}{\kw{then}}
\newcommand{\elsek}{\kw{else}}
\newcommand{\whilek}{\kw{while}}
\newcommand{\dok}{\kw{do}}
\newcommand{\prock}{\kw{proc}}
\newcommand{\sem}[1]{[\![#1]\!]}
\newcommand{\upd}[3]{#1[#2 \mapsto #3]}          % sigma[x -> n]
\newcommand{\bigstep}[3]{\langle #1, #2\rangle \Downarrow #3}
\newcommand{\triple}[3]{\{#1\}\, #2 \,\{#3\}}
\newcommand{\vdashH}{\vdash_{\mathrm{H}}}        % derivable in the proof system
\newcommand{\agree}[3]{#2 \sim_{#1} #3}          % agreement on a frame
\newcommand{\fv}{\operatorname{fv}}
              + TEXINPUTS=../common:
%    tex/note/main.tex   % ===========================================================================
%  Shared preamble for the Interproof demo.
%
%  It sits OUTSIDE either document's root, which is the point: a real project
%  keeps one notation file and both the paper and the notes read it, so the
%  build has to cope with a source file that is not under `[[document]].root`.
%  SyncTeX records it as an input, Interproof drops it (it holds no labels and
%  no page of its own), and everything keeps working.
%
%  The two documents reach it two different ways, on purpose:
%    tex/paper/main.tex  \input{preamble}              + TEXINPUTS=../common:
%    tex/note/main.tex   \input{../common/preamble}    plain relative path
%  See the comments in ../../interproof.toml.
% ===========================================================================

\usepackage[T1]{fontenc}
\usepackage[margin=1.1in]{geometry}
\usepackage{amsmath,amssymb,amsthm}
\usepackage{hyperref}
\usepackage{cleveref}          % MUST be loaded after hyperref

\hypersetup{colorlinks=true, linkcolor=[rgb]{0.10,0.25,0.55},
            citecolor=[rgb]{0.10,0.25,0.55}}

% ---- theorem environments -------------------------------------------------
% One shared counter, so `Definition 2.1` and `Lemma 2.2` interleave the way a
% reader expects.  The environment names here are exactly the ones listed in
% `[grammar].environments`; adding one means adding it in both places.
\theoremstyle{definition}
\newtheorem{definition}{Definition}[section]
\newtheorem{example}[definition]{Example}
\theoremstyle{plain}
\newtheorem{theorem}[definition]{Theorem}
\newtheorem{lemma}[definition]{Lemma}
\newtheorem{proposition}[definition]{Proposition}
\newtheorem{corollary}[definition]{Corollary}
\theoremstyle{remark}
\newtheorem{remark}[definition]{Remark}

\crefname{definition}{Definition}{Definitions}
\Crefname{definition}{Definition}{Definitions}
\crefname{lemma}{Lemma}{Lemmas}
\Crefname{lemma}{Lemma}{Lemmas}
\crefname{proposition}{Proposition}{Propositions}
\Crefname{proposition}{Proposition}{Propositions}
\crefname{corollary}{Corollary}{Corollaries}
\Crefname{corollary}{Corollary}{Corollaries}
\crefname{remark}{Remark}{Remarks}
\Crefname{remark}{Remark}{Remarks}

% ---- notation -------------------------------------------------------------
\newcommand{\Var}{\mathsf{Var}}
\newcommand{\Nat}{\mathbb{N}}
\newcommand{\Stores}{\Sigma}
\newcommand{\Assn}{\mathcal{A}}
\newcommand{\kw}[1]{\mathsf{#1}}
\newcommand{\skipk}{\kw{skip}}
\newcommand{\ifk}{\kw{if}}
\newcommand{\thenk}{\kw{then}}
\newcommand{\elsek}{\kw{else}}
\newcommand{\whilek}{\kw{while}}
\newcommand{\dok}{\kw{do}}
\newcommand{\prock}{\kw{proc}}
\newcommand{\sem}[1]{[\![#1]\!]}
\newcommand{\upd}[3]{#1[#2 \mapsto #3]}          % sigma[x -> n]
\newcommand{\bigstep}[3]{\langle #1, #2\rangle \Downarrow #3}
\newcommand{\triple}[3]{\{#1\}\, #2 \,\{#3\}}
\newcommand{\vdashH}{\vdash_{\mathrm{H}}}        % derivable in the proof system
\newcommand{\agree}[3]{#2 \sim_{#1} #3}          % agreement on a frame
\newcommand{\fv}{\operatorname{fv}}
    plain relative path
%  See the comments in ../../interproof.toml.
% ===========================================================================

\usepackage[T1]{fontenc}
\usepackage[margin=1.1in]{geometry}
\usepackage{amsmath,amssymb,amsthm}
\usepackage{hyperref}
\usepackage{cleveref}          % MUST be loaded after hyperref

\hypersetup{colorlinks=true, linkcolor=[rgb]{0.10,0.25,0.55},
            citecolor=[rgb]{0.10,0.25,0.55}}

% ---- theorem environments -------------------------------------------------
% One shared counter, so `Definition 2.1` and `Lemma 2.2` interleave the way a
% reader expects.  The environment names here are exactly the ones listed in
% `[grammar].environments`; adding one means adding it in both places.
\theoremstyle{definition}
\newtheorem{definition}{Definition}[section]
\newtheorem{example}[definition]{Example}
\theoremstyle{plain}
\newtheorem{theorem}[definition]{Theorem}
\newtheorem{lemma}[definition]{Lemma}
\newtheorem{proposition}[definition]{Proposition}
\newtheorem{corollary}[definition]{Corollary}
\theoremstyle{remark}
\newtheorem{remark}[definition]{Remark}

\crefname{definition}{Definition}{Definitions}
\Crefname{definition}{Definition}{Definitions}
\crefname{lemma}{Lemma}{Lemmas}
\Crefname{lemma}{Lemma}{Lemmas}
\crefname{proposition}{Proposition}{Propositions}
\Crefname{proposition}{Proposition}{Propositions}
\crefname{corollary}{Corollary}{Corollaries}
\Crefname{corollary}{Corollary}{Corollaries}
\crefname{remark}{Remark}{Remarks}
\Crefname{remark}{Remark}{Remarks}

% ---- notation -------------------------------------------------------------
\newcommand{\Var}{\mathsf{Var}}
\newcommand{\Nat}{\mathbb{N}}
\newcommand{\Stores}{\Sigma}
\newcommand{\Assn}{\mathcal{A}}
\newcommand{\kw}[1]{\mathsf{#1}}
\newcommand{\skipk}{\kw{skip}}
\newcommand{\ifk}{\kw{if}}
\newcommand{\thenk}{\kw{then}}
\newcommand{\elsek}{\kw{else}}
\newcommand{\whilek}{\kw{while}}
\newcommand{\dok}{\kw{do}}
\newcommand{\prock}{\kw{proc}}
\newcommand{\sem}[1]{[\![#1]\!]}
\newcommand{\upd}[3]{#1[#2 \mapsto #3]}          % sigma[x -> n]
\newcommand{\bigstep}[3]{\langle #1, #2\rangle \Downarrow #3}
\newcommand{\triple}[3]{\{#1\}\, #2 \,\{#3\}}
\newcommand{\vdashH}{\vdash_{\mathrm{H}}}        % derivable in the proof system
\newcommand{\agree}[3]{#2 \sim_{#1} #3}          % agreement on a frame
\newcommand{\fv}{\operatorname{fv}}
              + TEXINPUTS=../common:
%    tex/note/main.tex   % ===========================================================================
%  Shared preamble for the Interproof demo.
%
%  It sits OUTSIDE either document's root, which is the point: a real project
%  keeps one notation file and both the paper and the notes read it, so the
%  build has to cope with a source file that is not under `[[document]].root`.
%  SyncTeX records it as an input, Interproof drops it (it holds no labels and
%  no page of its own), and everything keeps working.
%
%  The two documents reach it two different ways, on purpose:
%    tex/paper/main.tex  % ===========================================================================
%  Shared preamble for the Interproof demo.
%
%  It sits OUTSIDE either document's root, which is the point: a real project
%  keeps one notation file and both the paper and the notes read it, so the
%  build has to cope with a source file that is not under `[[document]].root`.
%  SyncTeX records it as an input, Interproof drops it (it holds no labels and
%  no page of its own), and everything keeps working.
%
%  The two documents reach it two different ways, on purpose:
%    tex/paper/main.tex  \input{preamble}              + TEXINPUTS=../common:
%    tex/note/main.tex   \input{../common/preamble}    plain relative path
%  See the comments in ../../interproof.toml.
% ===========================================================================

\usepackage[T1]{fontenc}
\usepackage[margin=1.1in]{geometry}
\usepackage{amsmath,amssymb,amsthm}
\usepackage{hyperref}
\usepackage{cleveref}          % MUST be loaded after hyperref

\hypersetup{colorlinks=true, linkcolor=[rgb]{0.10,0.25,0.55},
            citecolor=[rgb]{0.10,0.25,0.55}}

% ---- theorem environments -------------------------------------------------
% One shared counter, so `Definition 2.1` and `Lemma 2.2` interleave the way a
% reader expects.  The environment names here are exactly the ones listed in
% `[grammar].environments`; adding one means adding it in both places.
\theoremstyle{definition}
\newtheorem{definition}{Definition}[section]
\newtheorem{example}[definition]{Example}
\theoremstyle{plain}
\newtheorem{theorem}[definition]{Theorem}
\newtheorem{lemma}[definition]{Lemma}
\newtheorem{proposition}[definition]{Proposition}
\newtheorem{corollary}[definition]{Corollary}
\theoremstyle{remark}
\newtheorem{remark}[definition]{Remark}

\crefname{definition}{Definition}{Definitions}
\Crefname{definition}{Definition}{Definitions}
\crefname{lemma}{Lemma}{Lemmas}
\Crefname{lemma}{Lemma}{Lemmas}
\crefname{proposition}{Proposition}{Propositions}
\Crefname{proposition}{Proposition}{Propositions}
\crefname{corollary}{Corollary}{Corollaries}
\Crefname{corollary}{Corollary}{Corollaries}
\crefname{remark}{Remark}{Remarks}
\Crefname{remark}{Remark}{Remarks}

% ---- notation -------------------------------------------------------------
\newcommand{\Var}{\mathsf{Var}}
\newcommand{\Nat}{\mathbb{N}}
\newcommand{\Stores}{\Sigma}
\newcommand{\Assn}{\mathcal{A}}
\newcommand{\kw}[1]{\mathsf{#1}}
\newcommand{\skipk}{\kw{skip}}
\newcommand{\ifk}{\kw{if}}
\newcommand{\thenk}{\kw{then}}
\newcommand{\elsek}{\kw{else}}
\newcommand{\whilek}{\kw{while}}
\newcommand{\dok}{\kw{do}}
\newcommand{\prock}{\kw{proc}}
\newcommand{\sem}[1]{[\![#1]\!]}
\newcommand{\upd}[3]{#1[#2 \mapsto #3]}          % sigma[x -> n]
\newcommand{\bigstep}[3]{\langle #1, #2\rangle \Downarrow #3}
\newcommand{\triple}[3]{\{#1\}\, #2 \,\{#3\}}
\newcommand{\vdashH}{\vdash_{\mathrm{H}}}        % derivable in the proof system
\newcommand{\agree}[3]{#2 \sim_{#1} #3}          % agreement on a frame
\newcommand{\fv}{\operatorname{fv}}
              + TEXINPUTS=../common:
%    tex/note/main.tex   % ===========================================================================
%  Shared preamble for the Interproof demo.
%
%  It sits OUTSIDE either document's root, which is the point: a real project
%  keeps one notation file and both the paper and the notes read it, so the
%  build has to cope with a source file that is not under `[[document]].root`.
%  SyncTeX records it as an input, Interproof drops it (it holds no labels and
%  no page of its own), and everything keeps working.
%
%  The two documents reach it two different ways, on purpose:
%    tex/paper/main.tex  \input{preamble}              + TEXINPUTS=../common:
%    tex/note/main.tex   \input{../common/preamble}    plain relative path
%  See the comments in ../../interproof.toml.
% ===========================================================================

\usepackage[T1]{fontenc}
\usepackage[margin=1.1in]{geometry}
\usepackage{amsmath,amssymb,amsthm}
\usepackage{hyperref}
\usepackage{cleveref}          % MUST be loaded after hyperref

\hypersetup{colorlinks=true, linkcolor=[rgb]{0.10,0.25,0.55},
            citecolor=[rgb]{0.10,0.25,0.55}}

% ---- theorem environments -------------------------------------------------
% One shared counter, so `Definition 2.1` and `Lemma 2.2` interleave the way a
% reader expects.  The environment names here are exactly the ones listed in
% `[grammar].environments`; adding one means adding it in both places.
\theoremstyle{definition}
\newtheorem{definition}{Definition}[section]
\newtheorem{example}[definition]{Example}
\theoremstyle{plain}
\newtheorem{theorem}[definition]{Theorem}
\newtheorem{lemma}[definition]{Lemma}
\newtheorem{proposition}[definition]{Proposition}
\newtheorem{corollary}[definition]{Corollary}
\theoremstyle{remark}
\newtheorem{remark}[definition]{Remark}

\crefname{definition}{Definition}{Definitions}
\Crefname{definition}{Definition}{Definitions}
\crefname{lemma}{Lemma}{Lemmas}
\Crefname{lemma}{Lemma}{Lemmas}
\crefname{proposition}{Proposition}{Propositions}
\Crefname{proposition}{Proposition}{Propositions}
\crefname{corollary}{Corollary}{Corollaries}
\Crefname{corollary}{Corollary}{Corollaries}
\crefname{remark}{Remark}{Remarks}
\Crefname{remark}{Remark}{Remarks}

% ---- notation -------------------------------------------------------------
\newcommand{\Var}{\mathsf{Var}}
\newcommand{\Nat}{\mathbb{N}}
\newcommand{\Stores}{\Sigma}
\newcommand{\Assn}{\mathcal{A}}
\newcommand{\kw}[1]{\mathsf{#1}}
\newcommand{\skipk}{\kw{skip}}
\newcommand{\ifk}{\kw{if}}
\newcommand{\thenk}{\kw{then}}
\newcommand{\elsek}{\kw{else}}
\newcommand{\whilek}{\kw{while}}
\newcommand{\dok}{\kw{do}}
\newcommand{\prock}{\kw{proc}}
\newcommand{\sem}[1]{[\![#1]\!]}
\newcommand{\upd}[3]{#1[#2 \mapsto #3]}          % sigma[x -> n]
\newcommand{\bigstep}[3]{\langle #1, #2\rangle \Downarrow #3}
\newcommand{\triple}[3]{\{#1\}\, #2 \,\{#3\}}
\newcommand{\vdashH}{\vdash_{\mathrm{H}}}        % derivable in the proof system
\newcommand{\agree}[3]{#2 \sim_{#1} #3}          % agreement on a frame
\newcommand{\fv}{\operatorname{fv}}
    plain relative path
%  See the comments in ../../interproof.toml.
% ===========================================================================

\usepackage[T1]{fontenc}
\usepackage[margin=1.1in]{geometry}
\usepackage{amsmath,amssymb,amsthm}
\usepackage{hyperref}
\usepackage{cleveref}          % MUST be loaded after hyperref

\hypersetup{colorlinks=true, linkcolor=[rgb]{0.10,0.25,0.55},
            citecolor=[rgb]{0.10,0.25,0.55}}

% ---- theorem environments -------------------------------------------------
% One shared counter, so `Definition 2.1` and `Lemma 2.2` interleave the way a
% reader expects.  The environment names here are exactly the ones listed in
% `[grammar].environments`; adding one means adding it in both places.
\theoremstyle{definition}
\newtheorem{definition}{Definition}[section]
\newtheorem{example}[definition]{Example}
\theoremstyle{plain}
\newtheorem{theorem}[definition]{Theorem}
\newtheorem{lemma}[definition]{Lemma}
\newtheorem{proposition}[definition]{Proposition}
\newtheorem{corollary}[definition]{Corollary}
\theoremstyle{remark}
\newtheorem{remark}[definition]{Remark}

\crefname{definition}{Definition}{Definitions}
\Crefname{definition}{Definition}{Definitions}
\crefname{lemma}{Lemma}{Lemmas}
\Crefname{lemma}{Lemma}{Lemmas}
\crefname{proposition}{Proposition}{Propositions}
\Crefname{proposition}{Proposition}{Propositions}
\crefname{corollary}{Corollary}{Corollaries}
\Crefname{corollary}{Corollary}{Corollaries}
\crefname{remark}{Remark}{Remarks}
\Crefname{remark}{Remark}{Remarks}

% ---- notation -------------------------------------------------------------
\newcommand{\Var}{\mathsf{Var}}
\newcommand{\Nat}{\mathbb{N}}
\newcommand{\Stores}{\Sigma}
\newcommand{\Assn}{\mathcal{A}}
\newcommand{\kw}[1]{\mathsf{#1}}
\newcommand{\skipk}{\kw{skip}}
\newcommand{\ifk}{\kw{if}}
\newcommand{\thenk}{\kw{then}}
\newcommand{\elsek}{\kw{else}}
\newcommand{\whilek}{\kw{while}}
\newcommand{\dok}{\kw{do}}
\newcommand{\prock}{\kw{proc}}
\newcommand{\sem}[1]{[\![#1]\!]}
\newcommand{\upd}[3]{#1[#2 \mapsto #3]}          % sigma[x -> n]
\newcommand{\bigstep}[3]{\langle #1, #2\rangle \Downarrow #3}
\newcommand{\triple}[3]{\{#1\}\, #2 \,\{#3\}}
\newcommand{\vdashH}{\vdash_{\mathrm{H}}}        % derivable in the proof system
\newcommand{\agree}[3]{#2 \sim_{#1} #3}          % agreement on a frame
\newcommand{\fv}{\operatorname{fv}}
    plain relative path
%  See the comments in ../../interproof.toml.
% ===========================================================================

\usepackage[T1]{fontenc}
\usepackage[margin=1.1in]{geometry}
\usepackage{amsmath,amssymb,amsthm}
\usepackage{hyperref}
\usepackage{cleveref}          % MUST be loaded after hyperref

\hypersetup{colorlinks=true, linkcolor=[rgb]{0.10,0.25,0.55},
            citecolor=[rgb]{0.10,0.25,0.55}}

% ---- theorem environments -------------------------------------------------
% One shared counter, so `Definition 2.1` and `Lemma 2.2` interleave the way a
% reader expects.  The environment names here are exactly the ones listed in
% `[grammar].environments`; adding one means adding it in both places.
\theoremstyle{definition}
\newtheorem{definition}{Definition}[section]
\newtheorem{example}[definition]{Example}
\theoremstyle{plain}
\newtheorem{theorem}[definition]{Theorem}
\newtheorem{lemma}[definition]{Lemma}
\newtheorem{proposition}[definition]{Proposition}
\newtheorem{corollary}[definition]{Corollary}
\theoremstyle{remark}
\newtheorem{remark}[definition]{Remark}

\crefname{definition}{Definition}{Definitions}
\Crefname{definition}{Definition}{Definitions}
\crefname{lemma}{Lemma}{Lemmas}
\Crefname{lemma}{Lemma}{Lemmas}
\crefname{proposition}{Proposition}{Propositions}
\Crefname{proposition}{Proposition}{Propositions}
\crefname{corollary}{Corollary}{Corollaries}
\Crefname{corollary}{Corollary}{Corollaries}
\crefname{remark}{Remark}{Remarks}
\Crefname{remark}{Remark}{Remarks}

% ---- notation -------------------------------------------------------------
\newcommand{\Var}{\mathsf{Var}}
\newcommand{\Nat}{\mathbb{N}}
\newcommand{\Stores}{\Sigma}
\newcommand{\Assn}{\mathcal{A}}
\newcommand{\kw}[1]{\mathsf{#1}}
\newcommand{\skipk}{\kw{skip}}
\newcommand{\ifk}{\kw{if}}
\newcommand{\thenk}{\kw{then}}
\newcommand{\elsek}{\kw{else}}
\newcommand{\whilek}{\kw{while}}
\newcommand{\dok}{\kw{do}}
\newcommand{\prock}{\kw{proc}}
\newcommand{\sem}[1]{[\![#1]\!]}
\newcommand{\upd}[3]{#1[#2 \mapsto #3]}          % sigma[x -> n]
\newcommand{\bigstep}[3]{\langle #1, #2\rangle \Downarrow #3}
\newcommand{\triple}[3]{\{#1\}\, #2 \,\{#3\}}
\newcommand{\vdashH}{\vdash_{\mathrm{H}}}        % derivable in the proof system
\newcommand{\agree}[3]{#2 \sim_{#1} #3}          % agreement on a frame
\newcommand{\fv}{\operatorname{fv}}
    plain relative path
%  See the comments in ../../interproof.toml.
% ===========================================================================

\usepackage[T1]{fontenc}
\usepackage[margin=1.1in]{geometry}
\usepackage{amsmath,amssymb,amsthm}
\usepackage{hyperref}
\usepackage{cleveref}          % MUST be loaded after hyperref

\hypersetup{colorlinks=true, linkcolor=[rgb]{0.10,0.25,0.55},
            citecolor=[rgb]{0.10,0.25,0.55}}

% ---- theorem environments -------------------------------------------------
% One shared counter, so `Definition 2.1` and `Lemma 2.2` interleave the way a
% reader expects.  The environment names here are exactly the ones listed in
% `[grammar].environments`; adding one means adding it in both places.
\theoremstyle{definition}
\newtheorem{definition}{Definition}[section]
\newtheorem{example}[definition]{Example}
\theoremstyle{plain}
\newtheorem{theorem}[definition]{Theorem}
\newtheorem{lemma}[definition]{Lemma}
\newtheorem{proposition}[definition]{Proposition}
\newtheorem{corollary}[definition]{Corollary}
\theoremstyle{remark}
\newtheorem{remark}[definition]{Remark}

\crefname{definition}{Definition}{Definitions}
\Crefname{definition}{Definition}{Definitions}
\crefname{lemma}{Lemma}{Lemmas}
\Crefname{lemma}{Lemma}{Lemmas}
\crefname{proposition}{Proposition}{Propositions}
\Crefname{proposition}{Proposition}{Propositions}
\crefname{corollary}{Corollary}{Corollaries}
\Crefname{corollary}{Corollary}{Corollaries}
\crefname{remark}{Remark}{Remarks}
\Crefname{remark}{Remark}{Remarks}

% ---- notation -------------------------------------------------------------
\newcommand{\Var}{\mathsf{Var}}
\newcommand{\Nat}{\mathbb{N}}
\newcommand{\Stores}{\Sigma}
\newcommand{\Assn}{\mathcal{A}}
\newcommand{\kw}[1]{\mathsf{#1}}
\newcommand{\skipk}{\kw{skip}}
\newcommand{\ifk}{\kw{if}}
\newcommand{\thenk}{\kw{then}}
\newcommand{\elsek}{\kw{else}}
\newcommand{\whilek}{\kw{while}}
\newcommand{\dok}{\kw{do}}
\newcommand{\prock}{\kw{proc}}
\newcommand{\sem}[1]{[\![#1]\!]}
\newcommand{\upd}[3]{#1[#2 \mapsto #3]}          % sigma[x -> n]
\newcommand{\bigstep}[3]{\langle #1, #2\rangle \Downarrow #3}
\newcommand{\triple}[3]{\{#1\}\, #2 \,\{#3\}}
\newcommand{\vdashH}{\vdash_{\mathrm{H}}}        % derivable in the proof system
\newcommand{\agree}[3]{#2 \sim_{#1} #3}          % agreement on a frame
\newcommand{\fv}{\operatorname{fv}}
